\item (a) El interior de un cilindro hueco se mantiene a una temperatura $T_b$ mientras que el exterior está a una temperatura menor $T_a$ (figura P7.45). La pared del cilindro tiene una conductividad térmica $k$. Si ignora los efectos de los extremos, demuestre que la rapidez de conducción de energía desde la superficie interior hasta la exterior en la dirección radial es
$$ \frac{dQ}{dt} = 2\pi Lk \left[ \frac{T_b - T_a}{\ln(b/a)} \right] $$
\textit{Sugerencia:} El gradiente de temperatura es $dT/dr$. Una corriente de energía radial pasa a través de un cilindro concéntrico de área $2\pi rL$. \\
(b) La sección de pasajeros de un jet comercial tiene la forma de un tubo cilíndrico con una longitud de 35.0 m y un radio interno de 2.50 m. Sus paredes están alineadas con un material aislante de 6.00 cm de espesor y con una conductividad térmica de $4.00 \times 10^{-5}\text{ cal/s}\cdot\text{cm}\cdot^\circ\text{C}$. Un calentador debe mantener la temperatura interior a $25.0^\circ\text{C}$ mientras que la temperatura exterior es de $-35.0^\circ\text{C}$. ¿Qué potencia debe proporcionarse al calentador?
